\subsubsection{Wichtige Winkelwerte}

Für häufig auftretende Winkel sind die Werte der trigonometrischen Funktionen
exakt bekannt:

\begin{center}
  \IfClozeTF
    {\renewcommand{\arraystretch}{2.05}}
    {\renewcommand{\arraystretch}{1.6}}
  \setlength{\tabcolsep}{4pt}
  \begin{tabularx}{\linewidth}{>{\raggedleft\arraybackslash}p{3.6cm}|*{5}{>{\centering\arraybackslash}X}}
    \toprule
    $\alpha$ im Gradmaß
      & $0^\circ$ & $30^\circ$ & $45^\circ$ & $60^\circ$ & $90^\circ$\\
    $\alpha$ im Bogenmaß
      & $0$ & $\frac{\pi}{6}$ & $\frac{\pi}{4}$
      & $\frac{\pi}{3}$ & $\frac{\pi}{2}$\\
    \midrule
    $\sin(\alpha)$
      & \Cloze{2.2cm}{$\frac{\sqrt{0}}{2}\mathrlap{\textcolor{gray}{\, = 0 }}$}
      & \Cloze{2.2cm}{$\frac{\sqrt{1}}{2}\mathrlap{\textcolor{gray}{\, = \frac12 }}$}
      & \Cloze{2.2cm}{$\frac{\sqrt{2}}{2}$}
      & \Cloze{2.2cm}{$\frac{\sqrt{3}}{2}$}
      & \Cloze{2.2cm}{$\frac{\sqrt{4}}{2}\mathrlap{\textcolor{gray}{\, = 1 }}$}\\
    $\cos(\alpha)$
      & \Cloze{2.2cm}{$\frac{\sqrt{4}}{2}\mathrlap{\textcolor{gray}{\, = 1 }}$}
      & \Cloze{2.2cm}{$\frac{\sqrt{3}}{2}$}
      & \Cloze{2.2cm}{$\frac{\sqrt{2}}{2}$}
      & \Cloze{2.2cm}{$\frac{\sqrt{1}}{2}\mathrlap{\textcolor{gray}{\, = \frac12 }}$}
      & \Cloze{2.2cm}{$\frac{\sqrt{0}}{2}\mathrlap{\textcolor{gray}{\, = 0 }}$}\\
    $\tan(\alpha)$
      & \Cloze{1cm}{$0$}
      & \Cloze{1cm}{$\frac{1}{\sqrt3}$}
      & \Cloze{1cm}{$1$}
      & \Cloze{1cm}{$\sqrt3$}
      & \textemdash\\
    \bottomrule
  \end{tabularx}
\end{center}
%

\IfClozeTF{}{\clearpage}
